% GENERATED FILE. Edit canonical lessons, then run npm run generate:books.
% canonical-source-hash: fnv1a64:ff63a69135564877
% canonical-lessons: IT-C31-che, IT-C31-penso-che, IT-C31-chi, IT-C31-che-cosa, IT-C31-che-buono

\chapter{Che}
\label{ch:it-che}
\hlchaptermodality{31}

\begin{chapteropening}
\textbf{By the end of this chapter:} \emph{I can join two clauses with che, say what I think, ask who and what, and exclaim.}

The chapter's last lesson spends what the book already owns: che from four lessons back and buono from Chapter 1, which had spent twenty-nine chapters welded to a time of day and was never once said about anything.
\end{chapteropening}

\section[che]{che — the word that joins two sentences into one}

\label{lesson:IT-C31-che}



\begin{quote}
A new chapter, and its first word is the one Italian uses most.

\begin{itemize}
\raggedright
  \item \emph{Recall:} say \emph{c'è / ci sono} — \textbf{R1}, one lesson back, before it decays
  \item \emph{Recall:} say \emph{non} — \textbf{R2}, the first real retrieval, five to fifteen lessons back
  \item \emph{Recall:} say \emph{il cuore} — \textbf{R3}, durable, twenty to sixty lessons back
  \item \emph{Recall:} say \emph{l'occhio} — \textbf{R3}, durable again, a second word from that distance
  \item \emph{Recall:} say \emph{uno, due, tre, quattro, cinque} — \textbf{R4}, recognition at distance, eighty lessons back or more
\end{itemize}
\end{quote}

\subsection*{You'll want to know: one word, joining}
\begin{quote}
\textbf{che} — \textbf{that, which, who, whom}.
\end{quote}

So far every Italian sentence in this book has been a single clause. \textbf{Che} staples two of them together, and it does not care what kind of thing it points back at:

\begin{quote}
\textbf{l'amico che parla italiano} — the friend \textbf{who} speaks Italian \textbf{il caffè che prendo} — the coffee \textbf{that} I'm having \textbf{la persona che aspetto} — the person \textbf{whom} I'm waiting for
\end{quote}

English picks a different word each time — \emph{who} for people, \emph{that} or \emph{which} for things, \emph{whom} when it is the object. Italian uses \textbf{che} for all of them, never changes it, and never leaves it out. English can say \emph{the coffee I'm having}; Italian cannot drop the \emph{che}.

One thing worth noticing in those three examples: \emph{che} is the \textbf{subject} of \emph{parla} in the first and the \textbf{object} of \emph{prendo} in the second and third. The word is the same either way. You do not have to decide which it is.

\begin{cousinweb}
\textbf{Che} is what is left of Latin \textbf{quid} and \textbf{quem} after they fell together. Latin distinguished them by case; Italian gave that up and kept one invariable word.

The spelling has been explained to you twice already without \emph{che} itself appearing. The \textbf{piacere} lesson taught that \emph{c} goes soft before \emph{e} and \emph{i}. The \textbf{mi chiamo} lesson then said that the \emph{h} in \emph{chi} \textquotedblleft{}is there precisely to keep it hard.\textquotedblright{} That is exactly what the \emph{h} is doing here: \textbf{che} is \emph{keh}, not \emph{cheh}, because the \emph{h} blocks the softening. You have had that explanation since your third chapter, and this is the word it was for.
\end{cousinweb}

\subsection*{Your turn}
Say these aloud:

\begin{itemize}
\raggedright
  \item \emph{che}
  \item \emph{l'amico che parla italiano}
  \item \emph{il caffè che prendo}
  \item \emph{la persona che aspetto}
\end{itemize}

\begin{tcolorbox}[breakable,colback=teal!4,colframe=teal!35!black,title={Before you move on}]
Which English words does \emph{che} cover? (\emph{\textbf{Who, whom, that, which}} — all of them.) Does it change? (\textbf{No} — never.) Why is it said \emph{keh}? (\textbf{The \emph{h} keeps the \emph{c} hard before \emph{e}} — the \emph{mi chiamo} rule.)
\end{tcolorbox}
\section[penso che]{penso che — having an opinion out loud}

\label{lesson:IT-C31-penso-che}



\begin{quote}
One lesson back, and the word that unlocks this one.

\begin{itemize}
\raggedright
  \item \emph{Recall:} say \emph{che} — \textbf{R1}, one lesson back, before it decays
  \item \emph{Recall:} say \emph{non capisco} — \textbf{R2}, the first real retrieval, five to fifteen lessons back
  \item \emph{Recall:} say \emph{l'orecchio} — \textbf{R3}, durable, twenty to sixty lessons back
  \item \emph{Recall:} say \emph{la bocca} — \textbf{R3}, durable again, a second word from that distance
  \item \emph{Recall:} say \emph{sei, sette, otto, nove, dieci} — \textbf{R4}, recognition at distance, eighty lessons back or more
\end{itemize}
\end{quote}

\subsection*{You'll want to know: the opinion frame}
The \emph{pensare} lesson gave you the verb \textquotedblleft{}to think\textquotedblright{} and its \emph{io} form \textbf{penso}. It gave you no way to say what you think. Here it is:

\begin{quote}
\textbf{Penso che …} — I think that …
\end{quote}

\begin{quote}
\textbf{Penso che Marco parla italiano.} — I think Marco speaks Italian. \textbf{Penso che il caffè è buono.} — I think the coffee is good. \textbf{Non penso che è vero.} — I don't think it's true.
\end{quote}

The \textbf{che} is the same invariable word as one lesson ago, doing its other job: there it hooked a clause onto a noun, and here it hooks a whole clause onto a verb. Italian will not let you leave it out. English drops it freely — \emph{I think he speaks Italian} — and Italian never does.

\textbf{A note for later, said now so it is not a surprise.} Careful written Italian puts the verb after \emph{penso che} into a different set of endings called the \emph{congiuntivo}: \emph{penso che Marco \textbf{parli} italiano}. Spoken Italian uses the plain present constantly, and everything above is ordinary speech. You will be understood. The \emph{congiuntivo} is a whole chapter of its own and it is not this one.

\subsection*{Your turn}
Say these aloud:

\begin{itemize}
\raggedright
  \item \emph{Penso che …}
  \item \emph{Penso che il caffè è buono.}
  \item \emph{Non penso che è vero.}
\end{itemize}

\begin{tcolorbox}[breakable,colback=teal!4,colframe=teal!35!black,title={Before you move on}]
Say \textquotedblleft{}I think that Marco speaks Italian.\textquotedblright{} (\emph{\textbf{Penso che Marco parla italiano}}.) May you drop the \emph{che}, as English drops \emph{that}? (\textbf{No}.) Which chapter gave you \emph{penso}? (\textbf{The} \emph{\textbf{pensare}} \textbf{lesson}.)
\end{tcolorbox}
\section[chi]{chi — who?}

\label{lesson:IT-C31-chi}



\begin{quote}
Three lessons of the chapter behind you.

\begin{itemize}
\raggedright
  \item \emph{Recall:} say \emph{penso che} — \textbf{R1}, one lesson back, before it decays
  \item \emph{Recall:} say \emph{né … né} — \textbf{R2}, the first real retrieval, five to fifteen lessons back
  \item \emph{Recall:} say \emph{il naso} — \textbf{R3}, durable, twenty to sixty lessons back
  \item \emph{Recall:} say \emph{lo stomaco} — \textbf{R3}, durable again, a second word from that distance
  \item \emph{Recall:} say \emph{lunedì, martedì, mercoledì, giovedì, venerdì} — \textbf{R4}, recognition at distance, eighty lessons back or more
\end{itemize}
\end{quote}

\subsection*{You'll want to know: The word}
\begin{quote}
\textbf{chi?} — \textbf{who?}
\end{quote}

\begin{quote}
\textbf{Chi è?} · \textbf{Chi parla italiano?} · \textbf{Chi aspetti?}
\end{quote}

Until this line, this book could ask \textbf{how} somebody was, \textbf{what} they were called and \textbf{how many} years they had, and it could not ask \textbf{who} anybody was. \textbf{Chi è?} is also what an Italian says through a closed door.

Keep \emph{chi} and \emph{che} apart by their jobs: \textbf{chi} asks about people and only about people. \textbf{Che} joins clauses, and — next lesson — asks about things.

\begin{cousinweb}
\textbf{Chi} is Latin \textbf{quī}, \textquotedblleft{}who\textquotedblright{}. The \textbf{chiedere} lesson taught you the sound change that produced it: Latin \emph{qu-} before a stressed front vowel went to Italian \textbf{chi-}. That lesson used it on \emph{quaerere} $\to$ \emph{chiedere}. It works the same way here.

So the \emph{ch} is hard for the same reason it is hard in \emph{che} and in \emph{mi chiamo}, and the word is \emph{kee}.
\end{cousinweb}

\subsection*{Your turn}
Say these aloud:

\begin{itemize}
\raggedright
  \item \emph{chi?}
  \item \emph{Chi è?}
  \item \emph{Chi parla italiano?}
\end{itemize}

\begin{tcolorbox}[breakable,colback=teal!4,colframe=teal!35!black,title={Before you move on}]
How do you ask \textquotedblleft{}who?\textquotedblright{} (\emph{\textbf{Chi}}.) Which lesson already taught you the \emph{qu-} $\to$ \emph{chi-} change? (\textbf{The} \emph{\textbf{chiedere}} \textbf{lesson}.)
\end{tcolorbox}
\section[che cosa]{che cosa — what?}

\label{lesson:IT-C31-che-cosa}



\begin{quote}
The fourth of the chapter's five, with a long throw back to the weekdays.

\begin{itemize}
\raggedright
  \item \emph{Recall:} say \emph{chi} — \textbf{R1}, one lesson back, before it decays
  \item \emph{Recall:} say \emph{anche / neanche} — \textbf{R2}, the first real retrieval, five to fifteen lessons back
  \item \emph{Recall:} say \emph{la gola} — \textbf{R3}, durable, twenty to sixty lessons back
  \item \emph{Recall:} say \emph{primo / prima} — \textbf{R3}, durable again, a second word from that distance
  \item \emph{Recall:} say \emph{sabato, domenica} — \textbf{R4}, recognition at distance, eighty lessons back or more
\end{itemize}
\end{quote}

\subsection*{You'll want to know: three ways to ask the same question}
\begin{quote}
\textbf{che cosa?} — \textbf{what?}
\end{quote}

Italian will let you say it three ways, and all three are ordinary:

\begin{quote}
\textbf{Che cosa prendi} · \textbf{Cosa prendi} · \textbf{Che prendi}
\end{quote}

\emph{Che cosa} is the full form, \emph{cosa} is commonest in the north, \emph{che} in the centre and south. Learn \emph{che cosa} and you will be understood everywhere.

\begin{quote}
\textbf{Che cosa è?} — what is it, set against \emph{\textbf{Chi è?}}, who is it.
\end{quote}

And the third member of the question set now comes unwelded. The age lesson gave you \textbf{Quanti anni hai?} as one lump; the \textbf{quanti} in it is a word, \textbf{quanto}, \textquotedblleft{}how much, how many\textquotedblright{}, agreeing with what it counts — \emph{quanti anni}, \emph{quante persone}. It is not glued to ages.

So: \textbf{chi} for people, \textbf{che cosa} for things, \textbf{quanto} for amounts.

\begin{cousinweb}
\textbf{Cosa} is Latin \textbf{causa} — \textquotedblleft{}a legal case, a matter at issue.\textquotedblright{} Roman lawyers used it constantly for \emph{the matter in hand}, and it drifted from \emph{the matter} to \emph{any matter} to, at last, \emph{thing}.

Every western Romance language took that road: Italian \emph{cosa}, Spanish \emph{cosa}, French \emph{chose}, Portuguese \emph{coisa} — all of them a Roman courtroom word. The same \emph{causa} came into English unworn, through law, as \textbf{cause}.
\end{cousinweb}

\subsection*{Your turn}
Say these aloud:

\begin{itemize}
\raggedright
  \item \emph{che cosa}
  \item \emph{Che cosa prendi}
  \item \emph{quanti anni, quante persone}
\end{itemize}

\begin{tcolorbox}[breakable,colback=teal!4,colframe=teal!35!black,title={Before you move on}]
Give the word for \textquotedblleft{}what\textquotedblright{}. (\emph{\textbf{Che cosa}}.) Name the Latin word behind \emph{cosa}. (\emph{\textbf{Causa}}, a court case.) And gloss the \emph{quanti} of \emph{quanti anni}. (\textbf{How much, how many}.)
\end{tcolorbox}
\section[Che buono!]{Che buono! — the exclamation}

\label{lesson:IT-C31-che-buono}



\begin{quote}
The chapter's last lesson, and it introduces no new word at all.

\begin{itemize}
\raggedright
  \item \emph{Recall:} say \emph{che cosa} — \textbf{R1}, one lesson back, before it decays
  \item \emph{Recall:} say \emph{c'è / ci sono} — \textbf{R2}, the first real retrieval, five to fifteen lessons back
  \item \emph{Recall:} say \emph{secondo / seconda} — \textbf{R3}, durable, twenty to sixty lessons back
  \item \emph{Recall:} say \emph{terzo / terza} — \textbf{R3}, durable again, a second word from that distance
  \item \emph{Recall:} say \emph{ora} — \textbf{R4}, recognition at distance, eighty lessons back or more
\end{itemize}
\end{quote}

\subsection*{You'll want to know: che, for the third time}
Put \textbf{che} in front of a describing word and it stops joining and starts exclaiming:

\begin{quote}
\textbf{Che buono!} — How good!
\end{quote}

and with a noun instead:

\begin{quote}
\textbf{Che caffè!} — What a coffee!
\end{quote}

There is no verb, and none is wanted. It is the whole utterance.

\textbf{This lesson introduces no new word.} \emph{Che} was three lessons ago; \textbf{buono} has been yours since the first chapter. What that chapter never did was let you say \emph{buono} about anything: it lived inside \textbf{buongiorno}, \textbf{buonasera} and \textbf{buonanotte}, welded to a time of day.

You can also put a noun after it: \textbf{Che caffè!} — what a coffee!

\begin{grammarlens}[title={Grammar lens: an adjective after essere}]
The exclamation is the loud way. Here is the quiet one, and it is the ordinary Italian sentence pattern this book has gone its whole length without using:

\begin{quote}
\textbf{subject} + \textbf{essere} + \textbf{describing word}
\end{quote}

\begin{quote}
\textbf{Il caffè è buono.} · \textbf{Il vino è rosso.} · \textbf{L'acqua è buona.}
\end{quote}

Every piece of that was separately taught and never joined: \emph{essere}, \emph{buono}, the four colours, and the drinks. Until this line, \emph{il caffè è nero} was unsayable although both halves had been on the page for a hundred lessons.

The describing word agrees with what it describes, exactly as \emph{primo / prima} did: \textbf{il caffè è buono}, \textbf{l'acqua è buona}, \textbf{i caffè sono buoni}.
\end{grammarlens}

\subsection*{Your turn}
Say these aloud:

\begin{itemize}
\raggedright
  \item \emph{Che buono!}
  \item \emph{Il caffè è buono.}
  \item \emph{Il vino è rosso.}
  \item \emph{Che caffè!}
\end{itemize}

\begin{tcolorbox}[breakable,colback=teal!4,colframe=teal!35!black,title={Before you move on}]
How do you say \textquotedblleft{}how good!\textquotedblright{}? (\emph{\textbf{Che buono!}}) Which two words is it made of, and when did you get them? (\emph{\textbf{Che}}, \textbf{three lessons ago;} \emph{\textbf{buono}}, \textbf{the first chapter}.) Say \textquotedblleft{}the coffee is good\textquotedblright{} as a plain sentence. (\emph{\textbf{Il caffè è buono}}.)
\end{tcolorbox}
